% LTeX: language=en-US
\documentclass{mimosis}

\usepackage{metalogo} % unnecessary
\usepackage{lipsum}   % unnecessary

% This can be used to typeset bold font in table while still aligning the numbers correctly:
\usepackage{etoolbox}

% Hyperlinks
\definecolor{GoetheBlue}{rgb}{0,0.38,0.56}
\usepackage[%
  colorlinks = true,
  citecolor  = GoetheBlue,
  linkcolor  = GoetheBlue,
  urlcolor   = GoetheBlue,
  unicode,
  ]{hyperref}
\renewcommand\UrlFont{\normalfont\small}
\renewcommand\path[1]{{\normalfont\small\detokenize{#1}}}
\usepackage[nameinlink,capitalise]{cleveref}

\usepackage{bookmark}

%%%%%%%%%%%%%%%%%%%%%%%%%%%%%%%%%%%%%%%%%%%%%%%%%%%%%%%%%%%%%%%%%%%%%%%%
% Fonts
%%%%%%%%%%%%%%%%%%%%%%%%%%%%%%%%%%%%%%%%%%%%%%%%%%%%%%%%%%%%%%%%%%%%%%%%

% Set serif font:
\usepackage[lining,semibold,scaled=1.05]{ebgaramond}

% Set monospace font:
\usepackage[scale=0.85]{sourcecodepro}

% if amsthm is needed, must be loaded before newtxmath
\usepackage{amsthm}

% Set math font:
\usepackage[ebgaramond,vvarbb,subscriptcorrection]{newtxmath}

% load after all math to give access to bold math using \bm{..} command:
\usepackage{bm}

% UML
\usepackage{tikz}
\usepackage{tikz-uml}

% enumerate
\usepackage{enumitem}

% Acronym and glossary directories are rare in theoretical computer science.
% If possible, try to avoid them.
% \newacronym[description={Principal component analysis}]{PCA}{PCA}{principal component analysis}
% \newacronym                                            {SNF}{SNF}{Smith normal form}
% \newacronym[description={Topological data analysis}]   {TDA}{TDA}{topological data analysis}

% \newglossaryentry{LaTeX}{%
%   name        = {\LaTeX},
%   description = {A document preparation system},
%   sort        = {LaTeX},
% }

% \newglossaryentry{Real numbers}{%
%   name        = {$\real$},
%   description = {The set of real numbers},
%   sort        = {Real numbers},
% }

% \makeindex
% \makeglossaries

%%%%%%%%%%%%%%%%%%%%%%%%%%%%%%%%%%%%%%%%%%%%%%%%%%%%%%%%%%%%%%%%%%%%%%%%
% Ordinals
%%%%%%%%%%%%%%%%%%%%%%%%%%%%%%%%%%%%%%%%%%%%%%%%%%%%%%%%%%%%%%%%%%%%%%%%

\makeatletter
\@ifundefined{st}{%
  \newcommand{\st}{\textsuperscript{\textup{st}}\xspace}
}{}
\@ifundefined{rd}{%
  \newcommand{\rd}{\textsuperscript{\textup{rd}}\xspace}
}{}
\@ifundefined{nd}{%
  \newcommand{\nd}{\textsuperscript{\textup{nd}}\xspace}
}{}
\makeatother

\renewcommand{\th}{\textsuperscript{\textup{th}}\xspace}

%%%%%%%%%%%%%%%%%%%%%%%%%%%%%%%%%%%%%%%%%%%%%%%%%%%%%%%%%%%%%%%%%%%%%%%%
% Incipit
%%%%%%%%%%%%%%%%%%%%%%%%%%%%%%%%%%%%%%%%%%%%%%%%%%%%%%%%%%%%%%%%%%%%%%%%

\title{Leveraging Large Language Models to Support Functional Programming Education}
% \subtitle{TODO:Less important thesis subtitle}
\author{Ahmad Khdeir}
\date{\today}

\begin{document}

  \frontmatter
  \begin{titlepage}
  \vspace*{5cm}
  \makeatletter
  \begin{center}
    \begin{Huge}
      \@title
    \end{Huge}\\[0.1cm]
    %
    \begin{Large}
      \@subtitle
    \end{Large}\\
    %
    \emph{by}\\
    \@author\\[\baselineskip]
    \@date
    %
    \vfill
    A document submitted in partial fulfillment
    of the requirements for the degree of\\
    \emph{Bachelor of Science}\\
    at\\
    \textsc{Goethe University Frankfurt}\\[2\baselineskip]
    \includegraphics[width=2cm]{Sources/Goethe_kompakt}
  \end{center}
  \makeatother
\end{titlepage}

\newpage
\null
\thispagestyle{empty}
\newpage

  \include{Sources/Erklärung}
  % ------------------------
% Acknowledgments
% ------------------------
\begin{center}
  \textsc{Acknowledgments}
\end{center}
\addcontentsline{toc}{section}{Acknowledgments}
I would like to express my sincere gratitude to my supervisor Alperen Kantarcı for his guidance and support throughout this thesis. I also thank Prof.\ Dr.\ Visvanathan Ramesh for his insightful feedback. Special thanks to my friends and family for their continuous encouragement. This project would not have been possible without their help and motivation.
  \begin{center}
  \textsc{Abstract}
\end{center}
%
% Remove the following line and your abstract here:
TODO

  \tableofcontents

  \mainmatter
  %%%%%%%%%%%%%%%%%%%%%%%%%%%%%%%%%%%%%%%%%%%%%%%%%%%%%%%%%%%%%%%%%%%%%%%%
\chapter{Introduction}
%%%%%%%%%%%%%%%%%%%%%%%%%%%%%%%%%%%%%%%%%%%%%%%%%%%%%%%%%%%%%%%%%%%%%%%%

\section{Why Functional Programming?}

Many universities opt for functional programming as the introductory course for first-year computer science students due to several compelling reasons. Functional programming provides an excellent foundation in algorithms, as it allows algorithms to be expressed clearly and concisely, avoiding unnecessary syntax that can often complicate understanding and implementation~\cite{joosten1993teaching}. Another significant advantage of introducing functional programming early in the curriculum is its emphasis on algorithmic design rather than syntactic details. While imperative programming continues to dominate the field of computer science, teaching alternative programming paradigms such as functional programming is beneficial to broaden students' perspectives. Additionally, many professors and institutions favor functional programming for beginners because it naturally introduces students to essential areas like algebra, artificial intelligence, formal language theory, and specification, making it an effective starting point for these critical topics. Furthermore, it's important to acknowledge that Haskell is one example of a functional programming language.

%%%%%%%%%%%%%%%%%%%%%%%%%%%%%%%%%%%%%%%%%%%%%%%%%%%%%%%%%%%%%%%%%%%%%%%%
\section{Challenges with Learning Functional Programming}

Now, let's discuss some common problems students face while learning programming in general. Many students struggle with understanding abstract concepts, such as how loops function or how methods execute during program runtime~\cite{derus2012difficulties}. Another significant issue is students' difficulty comprehending programming structures and the process of building a complete program. Additionally, mastering the syntax of programming languages is another frequent challenge. Merely understanding pseudo-code is insufficient; students must also grasp specific syntax rules, which can become especially confusing when learning multiple languages simultaneously. Hence, beyond conceptual understanding, practical application and experimentation with code are critical.
Another difficulty students encounter is the lack of sufficient examples provided during lectures. Although this may be intentional, as instructors often encourage independent research, it can be particularly confusing for first-year students who might require more structured guidance. Additionally, available examples sometimes fail to address all possible cases discussed in class. Furthermore, students frequently face challenges with complex data structures~\cite{milne2002difficulties}. This complexity often prompts numerous questions that require immediate clarification. However, tutors or professors are rarely available at all times, and providing instantaneous, real-time answers is practically challenging.
%%%%%%%%%%%%%%%%%%%%%%%%%%%%%%%%%%%%%%%%%%%%%%%%%%%%%%%%%%%%%%%%%%%%%%%%

\section{Proposed Solutions}

Here comes Haskify—but what exactly is Haskify? How does it help, and what does it provide? Haskify is a web application featuring a Haskell code editor with syntax highlighting and an integrated real-time compiler. More importantly, it includes an AI assistant specifically trained for functional programming. Unlike typical general-purpose AI tools such as ChatGPT, Haskify's AI is highly specialized in Haskell and functional programming, providing real-time, around-the-clock support.
This AI assistant has been carefully trained to focus exclusively on functional programming topics, particularly Haskell, ensuring that unrelated queries are not addressed. Even for relevant queries, the AI emphasizes detailed explanations and guidance rather than providing direct solutions immediately. This approach mirrors the role of an experienced university tutor, promoting deeper learning and understanding.
Therefore, Haskify serves as an essential tool for anyone aiming to learn Haskell and functional programming effectively. Users can write code directly in the editor, where the AI assistant promptly identifies and explains errors—be they syntax, logical, or conceptual. Crucially, the assistant guides users through their mistakes by offering insightful hints and constructive explanations without simply handing over the answers, preserving the joy and fulfillment inherent in the learning process.


%%%%%%%%%%%%%%%%%%%%%%%%%%%%%%%%%%%%%%%%%%%%%%%%%%%%%%%%%%%%%%%%%%%%%%%%

  % \include{Sources/Preliminaries}
  % Add your own chapters here:
  % \include{Sources/mychapter}
  \include{Sources/Conclusion}
  %%%%%%%%%%%%%%%%%%%%%%%%%%%%%%%%%%%%%%%%%%%%%%%%%%%%%%%%%%%%%%%%%%%%%%%%
\chapter{System Design and Specification}
%%%%%%%%%%%%%%%%%%%%%%%%%%%%%%%%%%%%%%%%%%%%%%%%%%%%%%%%%%%%%%%%%%%%%%%%

%%%%%%%%%%%%%%%%%%%%%%%%%%%%%%%%%%%%%%%%%%%%%%%%%%%%%%%%%%%%%%%%%%%%%%%%
\section{Requirements Analysis}
%%%%%%%%%%%%%%%%%%%%%%%%%%%%%%%%%%%%%%%%%%%%%%%%%%%%%%%%%%%%%%%%%%%%%%%%

The goal of this section is to identify and define the needs and expectations of the users of the Haskify. Haskify is a web-based educational tool designed to support students learning functional programming through Haskell. It integrates an AI assistant, a Haskell code editor with a live compiler.
\\ 
\newline
\textbf{Primary Users:}
\begin{itemize}
  \item \textbf{Students:} Interact with the system to write code, receive feedback, and learn.
  \item \textbf{Educators:} Could review usage or adapt exercises in future extensions.
\end{itemize}
\\
\textbf{System Objectives:}
\begin{itemize}
  \item Provide a web-based Haskell code editor with syntax highlighting and compilation.
  \item Deliver AI-driven guidance and explanation for functional programming exercises.
  \item Enhance the learning experience through immediate feedback.
\end{itemize}

%%%%%%%%%%%%%%%%%%%%%%%%%%%%%%%%%%%%%%%%%%%%%%%%%%%%%%%%%%%%%%%%%%%%%%%%
\section{Functional Requirements}
%%%%%%%%%%%%%%%%%%%%%%%%%%%%%%%%%%%%%%%%%%%%%%%%%%%%%%%%%%%%%%%%%%%%%%%%

\begin{enumerate}[label=FR\arabic*:]
  \item The system shall allow users to write and edit Haskell code in the code editor.
  \item The system shall compile and execute Haskell code and display the output or errors.
  \item The system shall provide an AI assistant that accepts user questions and code-related queries.
  \item The AI assistant shall respond with educational hints or explanations related to Haskell and functional programming.
  \item The system shall detect and highlight syntax or runtime errors in user code.
\end{enumerate}

%%%%%%%%%%%%%%%%%%%%%%%%%%%%%%%%%%%%%%%%%%%%%%%%%%%%%%%%%%%%%%%%%%%%%%%%
\section{Non-Functional Requirements}
%%%%%%%%%%%%%%%%%%%%%%%%%%%%%%%%%%%%%%%%%%%%%%%%%%%%%%%%%%%%%%%%%%%%%%%%

\begin{enumerate}[label=NFR\arabic*:]
  \item \textbf{Usability:} The interface should be intuitive and accessible to first-year students.
  \item \textbf{Performance:} The code compilation and AI responses must complete within 2 seconds under normal load.
  \item \textbf{Availability:} The system should be available 24/7 with minimal downtime.
  \item \textbf{Scalability:} The system should support at least 100 concurrent users during peak times.
  \item \textbf{Maintainability:} The codebase should be modular and follow good software engineering practices.
  \item \textbf{Portability:} The app must run in all major browsers (Chrome, Firefox, Safari).
\end{enumerate}
These requirements were considered throughout the development process. For example, usability was prioritized by selecting Monaco Editor, which provides students with a familiar coding environment similar to Visual Studio Code. Performance concerns were addressed by limiting AI context and optimizing the Haskell interpreter's execution environment. We avoided complex UI frameworks to reduce load times and improve responsiveness.
To meet availability and scalability goals, the system was designed using stateless APIs and could be deployed on scalable platforms like Vercel or AWS. Portability was ensured by rigorous testing across modern browsers. These constraints directly influenced design choices and technology selection.

%%%%%%%%%%%%%%%%%%%%%%%%%%%%%%%%%%%%%%%%%%%%%%%%%%%%%%%%%%%%%%%%%%%%%%%%
\newpage
\section{Use Case Diagram}
%%%%%%%%%%%%%%%%%%%%%%%%%%%%%%%%%%%%%%%%%%%%%%%%%%%%%%%%%%%%%%%%%%%%%%%%

Below is the use case diagram for Haskify

\begin{figure}[h!]
\centering
\begin{tikzpicture}

% Use-Case-System begin
\begin{umlsystem} [x=0, y=0, fill=pink!20] {{Haskify}}
\umlusecase[name=case1, width=2cm, x=1]{Write Haskell Code}
\umlusecase[name=case2, width=2cm, y=-7] {Run Code} 
\umlusecase[name=case3, width=2cm, y=-3.75] {Chat with AI Assistant} 
\umlusecase[name=case4, width=2cm, x=7, y=-3] {View Output / Errors} 
\umlusecase[name=case5, width=2.5cm, x=7, y=-6] {Receive Explanation / Feedback} 
\umlusecase[name=case6, width=2.5cm, x=7, y=-9] {return Explanation / Feedback} 
\umlusecase[name=case7, width=2cm, y=-9.5] {Return Output / Errors} 
\end{umlsystem}
% end Use-Case-System 
%Actors
\umlactor[x=-4, y=0]{Student}
\umlactor[x=-4, y=-4.5]{AI Assistant}
\umlactor[x=-4, y=-8]{Compiler Service}
%Relations
\umlassoc{Student}{case1}
\umlassoc{Student}{case2}
\umlassoc{Student}{case3}
\umlassoc{Student}{case4}
\umlassoc{Student}{case5}
\umlassoc{AI Assistant}{case6}
\umlassoc{Compiler Service}{case7}
\umlinclude{case2}{case7}
\umlinclude{case3}{case5}
\end{tikzpicture}
\end{figure}
\subsection*{\LARGE Use Case Specifications}
\textbf{Use Case: Write Haskell Code}
\begin{itemize}
  \item \textbf{Actor:} Student
  \item \textbf{Description:} The student writes Haskell code in the integrated editor.
  \item \textbf{Preconditions:} The user has opened the editor interface.
  \item \textbf{Postconditions:} The code is stored in memory and ready to be compiled.
\end{itemize}

\vspace{0.3cm}
\textbf{Use Case: Run Code}
\begin{itemize}
  \item \textbf{Actor:} Student
  \item \textbf{Description:} The user triggers code execution. The code is sent to the backend compiler.
  \item \textbf{Preconditions:} Valid Haskell code is present in the editor.
  \item \textbf{Postconditions:} Output or errors are returned and displayed.
\end{itemize}

\vspace{0.3cm}
\textbf{Use Case: Chat with AI Assistant}
\begin{itemize}
  \item \textbf{Actor:} Student
  \item \textbf{Description:} The student sends a query about Haskell code or concepts.
  \item \textbf{Preconditions:} An AI session is active.
  \item \textbf{Postconditions:} A helpful, domain-specific response is displayed.
\end{itemize}

\vspace{0.3cm}
\textbf{Use Case: Receive Explanation / Feedback}
\begin{itemize}
  \item \textbf{Actor:} AI Assistant
  \item \textbf{Description:} The assistant analyzes context and provides an explanation, hint, or guidance.
  \item \textbf{Preconditions:} A relevant user query was received.
  \item \textbf{Postconditions:} A response is generated and sent to the user.
\end{itemize}

\vspace{0.3cm}
\textbf{Use Case: Return Output / Errors}
\begin{itemize}
  \item \textbf{Actor:} Compiler Service
  \item \textbf{Description:} Executes code and sends output or error information.
  \item \textbf{Preconditions:} Valid compile request received.
  \item \textbf{Postconditions:} Result is returned to frontend for display.
\end{itemize}

%%%%%%%%%%%%%%%%%%%%%%%%%%%%%%%%%%%%%%%%%%%%%%%%%%%%%%%%%%%%%%%%%%%%%%%%
\newpage
\section{Technology Stack}
%%%%%%%%%%%%%%%%%%%%%%%%%%%%%%%%%%%%%%%%%%%%%%%%%%%%%%%%%%%%%%%%%%%%%%%%

Haskify is implemented using modern web technologies combined with cloud-based AI and compiler integration. The system is structured into a frontend and backend architecture to ensure scalability and maintainability.

\subsection{Frontend}

\begin{itemize}
  \item \textbf{React.js:} Used to build the user interface with component-based design.
  \item \textbf{Vite:} A fast development server and build tool for modern JavaScript applications.
  \item \textbf{HTML/CSS:} Used for layout and styling, including responsive design.
  \item \textbf{Monaco Editor:} A powerful, lightweight code editor (used in Visual Studio Code) embedded into the web application to support Haskell syntax, error highlighting, and real-time editing capabilities.
\end{itemize}

\subsection{Backend}

\begin{itemize}
  \item \textbf{Node.js:} JavaScript runtime used to build the backend server and API routes.
  \item \textbf{Express.js:} Backend web framework for handling API requests.
  \item \textbf{Nodemailer:} Handles email functionality (e.g., contact form).
  \item \textbf{Haskell (GHCi):} Used for compiling and running submitted Haskell code in a sandboxed environment.
\end{itemize}

\subsection{AI Assistant}

\begin{itemize}
  \item \textbf{OpenAI GPT API:} Powers the AI assistant with a specialized focus on functional programming guidance.
  \item \textbf{Custom prompt engineering:} Ensures AI answers remain educational and within domain scope.
\end{itemize}

\subsection{Other Tools and Services}

\begin{itemize}
  \item \textbf{GitHub Pages:} Used for deployment of the frontend.
  \item \textbf{GitHub:} Version control and collaborative code management.
\end{itemize}


  % This ensures that the subsequent sections are being included as root
  % items in the bookmark structure of your PDF reader.
  \bookmarksetup{startatroot}
  \backmatter

  % The acronyms and glossary are optional.
  % \begingroup
  %   \let\clearpage\relax
  %   \glsaddall
  %   \printglossary[type=\acronymtype]
  %   \newpage
  %   \printglossary
  % \endgroup

  % The index is optional
  % \printindex

  % The bibliography is required.
  \bibliographystyle{plain}% typically plainurl or alphaurl
  \bibliography{Thesis}
  

\end{document}
